\section{Verified Knowledge Compilation}\label{sec:prob}

Knowledge compilation is the route that makes a neural network and a probabilistic logic program differentiable end-to-end. A network's softmax becomes a distribution over weighted facts; provenance tracking yields a weighted Boolean formula; and that formula compiles into an arithmetic circuit whose forward pass is exact weighted model counting and whose backward pass yields exact gradients. Probabilistic systems such as ProbLog~\cite{problog2} follow this compile-once/evaluate-many model on the CPU. xlog instead keeps the compilation and evaluation data planes in CUDA buffers while the host orchestrates stages and reads bounded control state. Before use, xlog certifies the final smoothed circuit against its back-end source formula; this certificate does not cover provenance extraction or CNF encoding.

\subsection{The compilation pipeline}

The pipeline transforms a probabilistic program into a GPU-resident arithmetic circuit (the XGCF format) through five host-orchestrated stages (\Cref{fig:kc}). Their intermediate data stays in device buffers, while solver completion is observed through bounded status and error scalars.

\begin{figure*}[tbp]
  \centering
  \begin{tikzpicture}[x=1mm,y=1mm,
      stage/.style={fproc,anchor=north west,text width=24mm,
                    inner xsep=2mm,inner ysep=1.8mm,minimum height=14.5mm}]
    % ---- forward pipeline ----
    \node[stage] (s1) at (0,0)
      {provenance $\varphi$\\[0.4mm]\fsubt{hash-consing}};
    \node[stage] (s2) at (34,0)
      {CNF $\Phi$ (Tseitin)\\[0.4mm]\fsubt{prefix-sum var IDs}};
    \node[stage] (s3) at (68,0)
      {d-DNNF (D4)\\[0.4mm]\fsubt{BFS + block DFS}};
    \node[stage] (s4) at (102,0)
      {XGCF circuit $C$\\[0.4mm]\fsubt{levelized DAG}};
    \node[stage] (s5) at (142,0)
      {evaluate\\[0.4mm]\fsubt{forward $\log Z$}\\[-0.2mm]\fsubt{backward $\nabla$}};
    \node[diamond,aspect=1.2,draw=figaccent,fill=white,line width=0.8pt,
          inner sep=0.5pt] (gate) at (136,-7.25) {\scriptsize\checkmark};
    \draw[farrow] (s1) -- (s2);
    \draw[farrow] (s2) -- (s3);
    \draw[farrow] (s3) -- (s4);
    \draw[farrow] (s4) -- (gate);
    \draw[farrow] (gate) -- (s5);
    % ---- verification lane ----
    \node[fcert,anchor=north,text width=78mm,inner ysep=1.8mm] (cert) at (82,-26)
      {on-GPU CDCL equivalence certificate\\[0.4mm]
       $\varphi\wedge\lnot C$ and $C\wedge\lnot\varphi$ both \textsc{unsat}
       $\Rightarrow C\equiv\varphi$\\[0.4mm]
       \fsubt{resolution proofs checked on-device}\\[-0.2mm]
       \fsubt{complete proof, not sampling}};
    \draw[fver] (s1.south) |- (cert.west);
    \node[flabel,anchor=west] at (15,-20.5) {$\varphi$};
    \draw[fver] (s4.south) -- (s4.south |- cert.north);
    \node[flabel,anchor=west] at (117,-20.5) {$C$};
    \draw[fver] (cert.east) -| (gate.south);
    \node[flabel,anchor=west] at (137.5,-19) {proven $C\equiv\varphi$};
    \node[fnode,anchor=north,draw=figaccent,font=\scriptsize,
          text=figaccent,inner ysep=1.4mm,minimum height=5mm] (typederr) at (82,-49)
      {$\times$ typed verification error};
    \draw[fver] (cert.south) -- (typederr.north);
    \node[flabel,anchor=west] at (83.5,-45.8) {otherwise};
    % ---- GPU device plane ----
    \begin{scope}[on background layer]
      \node[fplane,fit=(s1)(s5)(cert)(typederr),inner sep=3mm] (plane) {};
    \end{scope}
    \node[fplanetitle,anchor=west] at ([xshift=4mm]plane.north west)
      {\planetitle{gpu-resident data plane}};
  \end{tikzpicture}
  \caption{Verified knowledge-compilation pipeline. The data plane uses CUDA
  kernels and device-resident state under host orchestration. The final smoothed
  circuit $C$ reaches evaluation only through the equivalence gate (\checkmark):
  an on-GPU CDCL solver proves $\varphi\wedge\lnot C$ and
  $C\wedge\lnot\varphi$ both unsatisfiable; the host observes bounded status and
  error scalars, and any failed or incomplete proof returns a typed error.}
  \label{fig:kc}
\end{figure*}

\textbf{Provenance extraction.} Provenance over deterministic evaluation produces a weighted Boolean formula~\cite{green2007provenance}. An on-device interner hash-conses node batches through an atomic GPU hash table, deduplicating structurally identical subformulas; this is essential because grounding can produce exponentially many redundant subformulas, and interning on the host would force the uninterned graph into host memory---reintroducing the very transfer the pipeline exists to avoid.

\textbf{CNF encoding.} A Tseitin~\cite{tseitin} transformation lowers the formula to CNF on the GPU, using parallel prefix sums to assign compact, gap-free variable IDs and emit clauses into a device-resident compressed-sparse-row structure.

\textbf{D4 compilation.} The CNF compiles into a Decision-DNNF (d-DNNF) circuit~\cite{darwiche2001dnnf} via a GPU-native variant of the D4 algorithm~\cite{d4compiler}: a BFS frontier exposes parallelism, per-block DFS workers perform component detection, decision-variable selection by VSIDS-like activity~\cite{chaff}, and recursive decomposition. A smoothing pass forces every branch of an OR/decision node to mention the same variable set, and the result is levelized so each topological level evaluates in one kernel launch. Decomposability, determinism, and smoothness (the structural properties that make weighted model counting over a Decision-DNNF correct~\cite{darwicheMarquis2002,chaviraDarwiche2008}) are established by construction. Certification runs after smoothing and therefore gates the exact final circuit that is cached and evaluated.

\subsection{Verification: a machine-checked certificate}

xlog does not trust the final compilation stage; it proves the final smoothed circuit $C$ equivalent to its back-end source formula $\varphi$ (\Cref{alg:verifiedkc}). It solves $\varphi \land \lnot C$ and $C \land \lnot \varphi$ on a GPU CDCL solver, and both must be unsatisfiable for $C \equiv \varphi$. UNSAT results carry resolution proofs checked on-device. The host reads bounded status and error scalars before accepting that GPU proof result; a wrong, incomplete, or budget-limited result returns a typed error and never reaches circuit caching or evaluation.

The certificate says that the circuit computes the same Boolean function as $\varphi$; the count-correctness invariants are guaranteed separately, by construction. The compilation back-end thus becomes a machine-checked stage instead of an assumption, closing the failure mode in which a silently miscompiled circuit would corrupt every downstream probability and gradient. The check covers the back-end stage specifically: provenance extraction and CNF encoding are trusted, not certified. CPU-side certified knowledge compilation exists~\cite{capelli2019}; the contribution here is performing the check on-device, sharing the GPU substrate. The same device-resident CDCL substrate is reused for SAT/MaxSAT queries and for epistemic candidate validation (\Cref{sec:epistemic}).

\begin{algorithm}[htbp]
\caption{Verified knowledge compilation. The final smoothed circuit is accepted only if proven equivalent to its back-end source formula.}
\label{alg:verifiedkc}
\begin{algorithmic}[1]
\Require weighted Boolean formula $\varphi$ (device-resident)
\Ensure verified circuit $C$ (XGCF), or typed error
\State $\varphi \gets \textsc{InternHashCons}(\varphi)$ \Comment{dedup via atomic GPU hash table}
\State $\Phi \gets \textsc{TseitinCNF}(\varphi)$ \Comment{prefix-sum variable IDs; CSR clauses}
\State $C \gets \textsc{D4}(\Phi)$; $C \gets \textsc{SmoothLevelize}(C)$
\State $u_1 \gets \textsc{CDCL}(\varphi \land \lnot C)$;\quad $u_2 \gets \textsc{CDCL}(C \land \lnot \varphi)$
\State observe bounded status and error scalars for $u_1,u_2$ on host
\If{$u_1 \neq \textsc{Unsat}$ \textbf{or} $u_2 \neq \textsc{Unsat}$}
  \State \Return typed verification error
\EndIf
\State check resolution proofs of $u_1, u_2$ on-device
\State \Return $C$ \Comment{cache the exact circuit that was certified}
\end{algorithmic}
\end{algorithm}

The verified circuit supports a level-parallel forward pass (log-space weighted model counting, one kernel per level, yielding $\log Z$) and a reverse-order backward pass that accumulates per-variable gradients. Both leave their results in GPU-resident buffers consumable directly by PyTorch via DLPack (\Cref{sec:nesy}).

\subsection{Circuit caching}

Building and certifying a circuit is the most expensive stage, on the order of a minute for a cold start on the MNIST-addition benchmark, and a naive implementation pays it every epoch. Measured separately, it is the equivalence proof and not the D4 compile that dominates that time (\Cref{sec:overhead}). The key observation is that the circuit \emph{structure} depends on the propositional structure of the grounded program and the evidence pattern, but \emph{not} on the numerical weights. When neural parameters change between iterations, the weights on probabilistic facts change while the Boolean formula and its compiled circuit stay identical. The certified final smoothed circuit is therefore cached, keyed by a structural hash of the provenance graph and evidence configuration; subsequent iterations update only the weight buffers and run the forward/backward kernels. On MNIST addition that amortization is what the circuit-cache ablation of \Cref{sec:cache-eval} measures.

\subsection{Monte Carlo alternative}

When the user selects approximate inference, xlog provides a GPU-parallel Monte Carlo engine for its supported fragment. It draws values for probabilistic facts on the GPU and evaluates the deterministic program in each sampled world, using rejection sampling or, when every evidence atom is a Bernoulli fact, evidence clamping that forces observed values and eliminates rejection waste. Both methods report confidence intervals. The resident sampler evaluates all worlds in a single device launch with a bounded device-side fixpoint. Its sampled core records zero tracked transfers; after synchronization, one pinned receipt returns status, counts, and trace metadata. That receipt's convergence status is authoritative: reaching the iteration bound returns \texttt{CONVERGENCE\_FAILURE}, even if partial counters exist. Fixed-count exhaustion returns \texttt{CAPACITY\_EXCEEDED}, byte-budget exhaustion returns \texttt{RESOURCE\_EXHAUSTED}, and unsupported fragments return their typed diagnostic; none selects a host execution path.
